\chapter{実験・評価}
% 【この章の目的】主張を検証できる実験設計と結果を示す．
% 内容入り例：examples/chapters/05-experiments-example.tex
% 結果の前に，データ分割，評価指標，比較手法，実装条件を固定する．

\section{データと評価手順}
% データ数，除外基準，分割単位，乱数seed，前処理，倫理・利用条件を確認する．

\section{比較条件と評価指標}
% ベースラインの選択理由，公平な条件，指標の定義と集計方法を示す．

\section{結果}
% 表・図は本文から必ず参照し，最良値だけでなくばらつきや失敗例も示す．表題は表の上，図題は図の下へ置く．

\section{考察}
% 観測事実，その解釈，限界，外的妥当性を分離して記述する．
